\section{Scaling 的新思考}

\subsection{语言如何塑造思维}

Ilya 提出了一个深刻的观察：\textbf{语言影响思维}。"Scaling"这一个词就足以驱动整个行业的行为——人们说"let's scale"，然后所有人都在做同样的事情。

\begin{importantbox}{Pre-training 的双重遗产}
两个术语深刻塑造了整个领域的思维：
\begin{enumerate}
    \item \textbf{AGI}：这个词源于对"narrow AI"（下棋 AI、游戏 AI）的反弹——既然窄 AI 不够好，那我们需要"general AI"
    \item \textbf{Pre-training}：它的 recipe 是"more pre-training $\rightarrow$ better at everything, uniformly"
\end{enumerate}
两者结合产生了"pre-training gives AGI"的思维定式。但问题是：\textbf{人类自身并不是 AGI}——人类缺乏大量知识，依靠的是 continual learning。
\end{importantbox}

\subsection{Research 时代不需要最大规模计算}

\begin{knowledgebox}{历史上的关键突破所需的计算量}
\begin{itemize}
    \item AlexNet：2 个 GPU
    \item Transformer 论文：最多 64 个 GPU（2017 年的，约等于今天 2 个 GPU）
    \item ResNet：类似的小规模
    \item O1 style reasoning：也不是世界上最重计算的东西
\end{itemize}
在研究阶段，你不需要绝对最大的计算规模来证明一个想法是正确的。大计算更多是在所有人都在同一个 paradigm 内竞争时的差异化因素。
\end{knowledgebox}

\subsection{SSI 的计算定位}

Ilya 解释了 SSI 的\$30 亿融资为何在研究中并不算少：
\begin{itemize}
    \item 大公司的大部分计算用于 inference（服务产品）
    \item 大公司需要大量员工（工程师、销售）和产品相关研究
    \item 真正用于纯研究的计算量，差距比表面数字小得多
\end{itemize}

\subsection{本章小结}

Scaling 时代"吸走了房间里所有的空气"，导致公司多于想法。回到 research 时代意味着重新重视创意和实验，而不仅仅是规模。关键突破不一定需要最大规模计算。


